% =========================================================
%  SECCIÓN IV — RESULTADOS
% =========================================================
\section{Resultados}

\subsection{Rendimiento del Kernel}

La Tabla~\ref{tab:rendimiento} resume los resultados cuantitativos
obtenidos en la medici\'on directa del tiempo de ejecuci\'on del
kernel para las cuatro versiones.
Los tiempos se midieron desde el propio programa C++ usando
\texttt{std::chrono::high\_resolution\_clock} sobre las
$N_{\text{tot}} = 52{,}428{,}800$ operaciones.
El rendimiento en GFLOPS se calcula como
$4 \times N_{\text{tot}} / t_{\text{kernel}}$ y el ancho de
banda utilizado como $3 \times 4 \times N_{\text{tot}} /
t_{\text{kernel}}$.

\begin{table}[htb]
\centering
\caption{Rendimiento de las cuatro versiones del kernel Prewitt.}
\label{tab:rendimiento}
\footnotesize
\begin{tabular}{@{}lcccc@{}}
\toprule
\textbf{Versi\'on} & \textbf{Tiempo (s)} & \textbf{GFLOPS} &
\textbf{BW (GB/s)} & \textbf{Speedup} \\
\midrule
Escalar     & 0.0354 & 5.92 & 17.76 & $1.00\times$ \\
Autom\'atica & 0.0286 & 7.34 & 22.02 & $\mathbf{1.24\times}$ \\
Guiada      & 0.0351 & 5.97 & 17.90 & $1.01\times$ \\
Expl\'icita  & 0.0293 & 7.15 & 21.44 & $1.21\times$ \\
\bottomrule
\end{tabular}
\end{table}

El speedup m\'aximo observado es de $1.24\times$ para la
vectorizaci\'on autom\'atica, seguido de $1.21\times$ para la
expl\'icita.
La versi\'on guiada ofrece pr\'acticamente el mismo rendimiento
que la escalar ($1.01\times$), a pesar de estar vectorizada con
AVX2.
Esto anticipa el comportamiento \textit{memory-bound} que se
confirmar\'a en el an\'alisis Roofline.

\subsection{An\'alisis Survey — Bucles Vectorizados}

Las Figuras~\ref{fig:loops-escalar} a~\ref{fig:loops-explicita}
muestran la lista de bucles identificados por Intel Advisor
para cada versi\'on, con sus tiempos de ejecuci\'on y tipo
de vectorizaci\'on detectado.

\begin{figure}[htb]
\begin{center}
\includegraphics[width=\columnwidth]{timeLoops-Escalar}
\end{center}
\caption{Survey de bucles --- versi\'on Escalar. El kernel
(\texttt{kernel\_escalar:107}) reporta tipo \textit{Scalar}
con ISA SSE (uso escalar de registros XMM, sin SIMD empaquetado),
conforme al flag \texttt{-no-vec}. Trip count: 6{,}553{,}600.}
\label{fig:loops-escalar}
\end{figure}

\begin{figure}[htb]
\begin{center}
\includegraphics[width=\columnwidth]{timeLoops-Auto}
\end{center}
\caption{Survey de bucles --- versi\'on Autom\'atica. El kernel
(\texttt{kernel\_auto:84}) se marca como \textit{Vectorized (Body)}
con ISA AVX2 y trip count de 3{,}276{,}800, lo que corresponde
a un ancho vectorial de $52{,}428{,}800 / 3{,}276{,}800 = 16$
floats por iteraci\'on.}
\label{fig:loops-auto}
\end{figure}

\begin{figure}[htb]
\begin{center}
\includegraphics[width=\columnwidth]{timeLoops-Guiada}
\end{center}
\caption{Survey de bucles --- versi\'on Guiada. El kernel
(\texttt{kernel\_guiado:90}) aparece como \textit{Vectorized (Body)}
con ISA AVX2 y trip count de 6{,}553{,}600, correspondiente a un
ancho vectorial de~8 floats, exactamente el valor especificado
con \texttt{simdlen(8)}.}
\label{fig:loops-guiada}
\end{figure}

\begin{figure}[htb]
\begin{center}
\includegraphics[width=\columnwidth]{timeLoops-Explicita}
\end{center}
\caption{Survey de bucles --- versi\'on Expl\'icita. El kernel
(\texttt{kernel\_explicito:93}) se reporta como \textit{Vectorized (Body)}
con ISA AVX2 y trip count de 1{,}638{,}400, lo que implica
un ancho efectivo de $52{,}428{,}800 / 1{,}638{,}400 = 32$
floats: el compilador aplic\'o adicionalmente un desenrollado
(\textit{unroll}) de $4\times$ sobre el bucle AVX2 escrito a mano.}
\label{fig:loops-explicita}
\end{figure}

La Tabla~\ref{tab:survey} resume los datos del Survey para los
kernels principales.

\begin{table}[htb]
\centering
\caption{Datos del Survey de Intel Advisor para el kernel Prewitt.}
\label{tab:survey}
\footnotesize
\begin{tabular}{@{}lcccc@{}}
\toprule
\textbf{Versi\'on} & \textbf{Tipo} & \textbf{ISA} &
\textbf{Trip count} & \textbf{Ancho vec.} \\
\midrule
Escalar     & Scalar           & SSE  & 6{,}553{,}600 & 1  \\
Autom\'atica & Vectorized (Body)& AVX2 & 3{,}276{,}800 & 16 \\
Guiada      & Vectorized (Body)& AVX2 & 6{,}553{,}600 &  8 \\
Expl\'icita  & Vectorized (Body)& AVX2 & 1{,}638{,}400 & 32$^*$ \\
\bottomrule
\multicolumn{5}{@{}l@{}}{\scriptsize $^*$ Incluye unroll $4\times$ aplicado por el compilador.}
\end{tabular}
\end{table}

\subsection{An\'alisis Roofline}

Las Figuras~\ref{fig:roofline-escalar} a~\ref{fig:roofline-explicita}
presentan el modelo Roofline generado por Intel Advisor para cada
versi\'on.

\begin{figure}[htb]
\begin{center}
\includegraphics[width=\columnwidth]{graficaRoofline-Escalar}
\end{center}
\caption{Modelo Roofline --- versi\'on Escalar. El punto rojo
se ubica con AI $\approx 0.33$~FLOP/Byte y $\approx 6$--8~GFLOPS,
en la zona acotada entre los techos de ancho de banda L2/L3 y el
techo escalar (9.28~GFLOPS).}
\label{fig:roofline-escalar}
\end{figure}

\begin{figure}[htb]
\begin{center}
\includegraphics[width=\columnwidth]{graficaRoofline-Auto}
\end{center}
\caption{Modelo Roofline --- versi\'on Autom\'atica. El punto se
sit\'ua en la misma zona que el Escalar; la ganancia vectorial es
m\'inima porque la memoria no puede alimentar los registros AVX2
a la velocidad m\'axima.}
\label{fig:roofline-auto}
\end{figure}

\begin{figure}[htb]
\begin{center}
\includegraphics[width=\columnwidth]{graficaRoofline-Guiada}
\end{center}
\caption{Modelo Roofline --- versi\'on Guiada. Comportamiento
id\'entico a las versiones anteriores. El ancho vectorial de~8
(menor al de la versi\'on autom\'atica) resulta en rendimiento
similar al escalar.}
\label{fig:roofline-guiada}
\end{figure}

\begin{figure}[htb]
\begin{center}
\includegraphics[width=\columnwidth]{graficaRoofline-Explicita}
\end{center}
\caption{Modelo Roofline --- versi\'on Expl\'icita. A pesar del
ancho efectivo de~32 floats (con unroll), el punto permanece en la
misma zona \textit{memory-bound}, con rendimiento ligeramente
superior al escalar.}
\label{fig:roofline-explicita}
\end{figure}

En los cuatro casos, el punto de operaci\'on se ubica en la regi\'on
denominada ``\textit{Bound by compute and memory roofs}'', a la
izquierda del punto de cresta del Roofline.
La intensidad aritm\'etica medida por Advisor para el programa
completo es de $\mathcal{I}_{\text{prog}} = 0.100$~FLOP/Byte,
inferior al valor te\'orico del kernel ($0.333$~FLOP/Byte),
porque Advisor incluye todas las operaciones de memoria del
programa (I/O de imagen PGM, c\'alculo de gradientes Prewitt,
replicaci\'on de datos), que diluyen la intensidad global.

\subsection{M\'etricas Generales del Programa}

Las Figuras~\ref{fig:metrics-escalar}
y~\ref{fig:metrics-auto} muestran los Program Metrics
reportados por Intel Advisor para las versiones Escalar y
Autom\'atica como versiones representativas.

\begin{figure}[htb]
\begin{center}
\includegraphics[width=\columnwidth]{summaryProgramMetrics-Escalar}
\end{center}
\caption{Program Metrics --- versi\'on Escalar. Tiempo total
$0.28$~s; rendimiento $0.76$~GFLOPS; tiempo del bucle vectorizado
(SSE escalar) $0.03$~s ($11\%$) frente a $0.24$~s ($89\%$) de
c\'odigo escalar (I/O, gradientes, replicaci\'on). ISA del
programa: AVX2, AVX, SSE2, SSE.}
\label{fig:metrics-escalar}
\end{figure}

\begin{figure}[htb]
\begin{center}
\includegraphics[width=\columnwidth]{summaryProgramMetrics-Auto}
\end{center}
\caption{Program Metrics --- versi\'on Autom\'atica. Tiempo total
$0.27$~s; rendimiento $0.76$~GFLOPS; tiempo del bucle AVX2
$0.03$~s ($11\%$). ISA del programa: AVX2, AVX, SSE2 (sin SSE
puro). La distribuci\'on de tiempos refleja que el kernel
vectorizado es solo una peque\~na fracci\'on del programa total.}
\label{fig:metrics-auto}
\end{figure}

La Tabla~\ref{tab:advisor-metricas} consolida los datos de
Program Metrics para las cuatro versiones.

\begin{table}[htb]
\centering
\caption{Resumen de m\'etricas de Intel Advisor (Program Metrics).}
\label{tab:advisor-metricas}
\footnotesize
\begin{tabular}{@{}lcccc@{}}
\toprule
\textbf{Versi\'on} & \textbf{Elapsed (s)} & \textbf{GFLOPS} &
\textbf{AI (FLOP/B)} & \textbf{Kernel (\%)} \\
\midrule
Escalar      & 0.28 & 0.76 & 0.100 & 11.1 \\
Autom\'atica  & 0.27 & 0.76 & 0.100 & 11.1 \\
Guiada       & 0.24 & 0.86 & 0.100 & 12.5 \\
Expl\'icita   & 0.26 & 0.82 & 0.100 & 12.0 \\
\bottomrule
\end{tabular}
\end{table}

Las Figuras~\ref{fig:metrics-guiada} y~\ref{fig:metrics-explicita}
complementan el panorama para las versiones Guiada y Expl\'icita.

\begin{figure}[htb]
\begin{center}
\includegraphics[width=\columnwidth]{summaryProgramMetrics-Guiada}
\end{center}
\caption{Program Metrics --- versi\'on Guiada. Tiempo total
$0.24$~s; rendimiento $0.86$~GFLOPS; bucle AVX2 con
\texttt{simdlen(8)} corresponde al $12.5\%$ del tiempo. Pese a
que el ancho vectorial es inferior al de la versi\'on autom\'atica,
el tiempo total es el menor de las cuatro, por variaciones de
timing en el c\'odigo escalar circundante.}
\label{fig:metrics-guiada}
\end{figure}

\begin{figure}[htb]
\begin{center}
\includegraphics[width=\columnwidth]{summaryProgramMetrics-Explicita}
\end{center}
\caption{Program Metrics --- versi\'on Expl\'icita. Tiempo total
$0.26$~s; rendimiento $0.82$~GFLOPS; bucle AVX2 el $12\%$ del
tiempo. Dos bucles escalares adicionales de c\'alculo de gradientes
Prewitt (\texttt{calcular\_gradientes:57}, \texttt{:62}) son
visibles en el Survey de esta versi\'on; en las dem\'as no superan
el umbral de medici\'on de Advisor.}
\label{fig:metrics-explicita}
\end{figure}

Un hallazgo relevante es que el tiempo total del programa est\'a
dominado ($\approx 88$--$89\%$) por c\'odigo escalar: lectura de
la imagen PGM, c\'alculo de gradientes Prewitt con las m\'ascaras
3$\times$3, replicaci\'on de datos en memoria y escritura del
resultado.
El kernel vectorizado representa \'unicamente el $11$--$12\%$
del tiempo de ejecuci\'on total.
Esta distribuci\'on es independiente de la estrategia de
vectorizaci\'on empleada, ya que las partes escalares del
programa son las mismas en las cuatro versiones.

\subsection{An\'alisis de C\'odigo Activo (Code Analytics)}

El panel Code Analytics de Intel Advisor identifica las funciones
que mayor tiempo de c\'omputo consumen dentro de los bucles
vectorizados.
En las cuatro versiones, la funci\'on m\'as activa reportada es
\texttt{\_\_intel\_avx\_rep\_memset}, que es la rutina interna
del compilador para inicializaci\'on de memoria usando
instrucciones AVX.
Esto confirma que la mayor parte del tiempo vectorizado
corresponde a transferencias de datos en memoria y no a
operaciones aritm\'eticas, consistente con el car\'acter
\textit{memory-bound} revelado por el Roofline.

\begin{figure}[htb]
\begin{center}
\includegraphics[width=\columnwidth]{codeAnalytics-Escalar}
\end{center}
\caption{Code Analytics --- versi\'on Escalar. El tiempo total
registrado en el bucle es de $0.16$~s, con un GINTOPS (giga
operaciones enteras por segundo) de speedup $6.9\times$ respecto
al techo escalar de referencia. La funci\'on dominante es
\texttt{\_\_intel\_avx\_rep\_memset}, lo que confirma que el
tiempo se gasta principalmente en movimiento de datos.}
\label{fig:analytics-escalar}
\end{figure}

\begin{figure}[htb]
\begin{center}
\includegraphics[width=\columnwidth]{codeAnalytics-Explicita}
\end{center}
\caption{Code Analytics --- versi\'on Expl\'icita. Tiempo de
$0.15$~s con speedup GINTOPS de $5.8\times$. Comparando con
la versi\'on Escalar ($6.9\times$), la diferencia en el perfil
de funciones activas confirma que los intr\'insecos AVX2
modifican el patr\'on de acceso a memoria respecto al c\'odigo
escalar, aunque ambas versiones siguen siendo dominadas por la
transferencia de datos.}
\label{fig:analytics-explicita}
\end{figure}
